\subsection{Directional Derivatives and The Gradient Vector}

\content{

}{% presentation
\vspace{3in}
}{% nopresentation
}{% comments
Demonstrate what the directional derivative means geometricaly, by drawing a
diagram with both partial derivatives on it, then a directional derivative at
the same point. 
}

\content{
\begin{definition} The directional derivative of $f(x,y)$ at $(x_0,y_0)$ in the
direction of a unit vector $\vec{u} = \la a,b\ra$, is given by 
$$ D_{\vec{u}}f(x_0,y_0) = \lim_{h\rightarrow 0} 
\frac{f(x_0+ha,y_0+hb)-f(x_0,y_0)}{h}. $$
\end{definition}
}{% presentation
}{% nopresentation
}{% comments
}

\content{
\begin{theorem} If $f$ is a differentiable function of $x$ and $y$, then $f$ has
a directional derivative in the direction of any unit vector $\vec{u} =\la
a,b\ra$, and 
$$ D_{\vec{u}}f(x,y) = f_x(x,y)a+f_y(x,y)b. $$
\end{theorem}
}{% presentation
\vspace{4in}
}{% nopresentation
}{% comments
Prove this by 
}

\subsubsection{The Gradient Vector}

\content{
Notice that we could write the formula for the directional derivative
$D_{\vec{u}} f(x,y)$ as 
$$ \la f_x(x,y), f_y(x,y)\ra \cdot \vec{u}. $$
\begin{definition} We call the vector 
$$\la f_x(x,y), f_y(x,y)\ra$$ the gradient vector of $f$ at $(x,y)$.
It is denoted $\nabla f$.
\end{definition}
}{% presentation
}{% nopresentation
}{% comments
Note that this formula also holds for functions of three variables. 
}

\content{
\begin{example} If $f(x,y)$ is given by 
$$ f(x,y) = \sin(x)\cos(y), $$
find the directional derivative of $f$ at the point $(0,0)$ in the direction $\vec{u} = \la 1,1\ra$.
\end{example}
}{% presentation
\vspace{3in}
}{% nopresentation
}{% comments
First have to make $\vec{u}$ a unit vector!
}

\content{
\begin{example} If $f(x,y,z) = x\sin(yz)$, find the gradient of $f$ and the
directional derivative of $f$ at (1,3,0) in the direction of $\la 1,2,-1\ra$.
\end{example}
}{% presentation
\vspace{2in}
}{% nopresentation
}{% comments
}

\content{
\begin{theorem} Suppose that $f$ is a differentiable function of two or three
variables. The maximum value of the directional derivative is given in the
direction of the gradient of $f$, and it's value is $|\nabla f|$.
\end{theorem}
}{% presentation
\vspace{2in}
}{% nopresentation
}{% comments
Prove this by using the dot product formula for the directional derivative, and
the fact that 
$$ \vec{u}\cdot \vec{v} = |\vec{u}||\vec{v}|\cos \theta. $$
Then since $\vec{u}$ will be a unit vector, we simply must maximize
$\cos\theta$.
}

\subsubsection{Tangent Planes to Level Surfaces}
\content{
\begin{definition} Given a surface defined by 
$$ F(x,y,z) = k, $$
the tangent plane to the surface at the point $(x_0,y_0,z_0)$ is given by 
$$ \pd{F}{x}(x_0,y_0,z_0)(x-x_0) + \pd{F}{y}(x_0,y_0,z_0)(y-y_0) + 
\pd{F}{z}(x_0,y_0,z_0)(z-z_0) = 0, $$
\end{definition}
}{% presentation
\vspace{2in}
}{% nopresentation
}{% comments
Sketch a quick picture of a  surface (level curve) and demonstrate that the
gradient is perpendicular to the surface at the given point, so the tangent
plane is simply the plane whose normal vector is the gradient. 
}

\content{
\begin{example} Find the tangent plane to the surface defined by 
$$ \frac{x^2}{4} + y^2 + \frac{z^2}{9}=3 $$
at the point $(-2,1,-3)$.
\end{example}
}{% presentation
\vspace{3in}
}{% nopresentation
}{% comments
Final plane is given by 
$$ -1(x+2)+2(y-1) -\frac{2}{3}(z+3) = 0. $$
}
